% Options for packages loaded elsewhere
\PassOptionsToPackage{unicode}{hyperref}
\PassOptionsToPackage{hyphens}{url}
\PassOptionsToPackage{dvipsnames,svgnames,x11names}{xcolor}
%
\documentclass[
  letterpaper,
  DIV=11,
  numbers=noendperiod]{scrartcl}

\usepackage{amsmath,amssymb}
\usepackage{setspace}
\usepackage{iftex}
\ifPDFTeX
  \usepackage[T1]{fontenc}
  \usepackage[utf8]{inputenc}
  \usepackage{textcomp} % provide euro and other symbols
\else % if luatex or xetex
  \usepackage{unicode-math}
  \defaultfontfeatures{Scale=MatchLowercase}
  \defaultfontfeatures[\rmfamily]{Ligatures=TeX,Scale=1}
\fi
\usepackage{lmodern}
\ifPDFTeX\else  
    % xetex/luatex font selection
\fi
% Use upquote if available, for straight quotes in verbatim environments
\IfFileExists{upquote.sty}{\usepackage{upquote}}{}
\IfFileExists{microtype.sty}{% use microtype if available
  \usepackage[]{microtype}
  \UseMicrotypeSet[protrusion]{basicmath} % disable protrusion for tt fonts
}{}
\usepackage{xcolor}
\setlength{\emergencystretch}{3em} % prevent overfull lines
\setcounter{secnumdepth}{5}
% Make \paragraph and \subparagraph free-standing
\ifx\paragraph\undefined\else
  \let\oldparagraph\paragraph
  \renewcommand{\paragraph}[1]{\oldparagraph{#1}\mbox{}}
\fi
\ifx\subparagraph\undefined\else
  \let\oldsubparagraph\subparagraph
  \renewcommand{\subparagraph}[1]{\oldsubparagraph{#1}\mbox{}}
\fi


\providecommand{\tightlist}{%
  \setlength{\itemsep}{0pt}\setlength{\parskip}{0pt}}\usepackage{longtable,booktabs,array}
\usepackage{calc} % for calculating minipage widths
% Correct order of tables after \paragraph or \subparagraph
\usepackage{etoolbox}
\makeatletter
\patchcmd\longtable{\par}{\if@noskipsec\mbox{}\fi\par}{}{}
\makeatother
% Allow footnotes in longtable head/foot
\IfFileExists{footnotehyper.sty}{\usepackage{footnotehyper}}{\usepackage{footnote}}
\makesavenoteenv{longtable}
\usepackage{graphicx}
\makeatletter
\def\maxwidth{\ifdim\Gin@nat@width>\linewidth\linewidth\else\Gin@nat@width\fi}
\def\maxheight{\ifdim\Gin@nat@height>\textheight\textheight\else\Gin@nat@height\fi}
\makeatother
% Scale images if necessary, so that they will not overflow the page
% margins by default, and it is still possible to overwrite the defaults
% using explicit options in \includegraphics[width, height, ...]{}
\setkeys{Gin}{width=\maxwidth,height=\maxheight,keepaspectratio}
% Set default figure placement to htbp
\makeatletter
\def\fps@figure{htbp}
\makeatother

% \usepackage{hyperref}
% \hypersetup{
%     colorlinks,
%     linkcolor={blue!100!black},
%     citecolor={blue!100!black},
%     urlcolor={blue!100!black}
% }
% \usepackage{apacite}
% \usepackage[round]{natbib} 
% \usepackage{graphicx}
% \usepackage{float}
% \usepackage{caption}
% \usepackage[toc,page]{appendix}
% \usepackage{booktabs,caption}
% \usepackage[flushleft]{threeparttable}
% \usepackage{tabularx}
\usepackage{sfmath}
\usepackage{fontspec}
\usepackage[utf8]{inputenc}
\usepackage{siunitx}
\usepackage{amsfonts}
\usepackage{amsmath}
% \usepackage{xcolor}
% \usepackage{scrextend}
% \deffootnote[2em]{2em}{1em}{\textsuperscript{\thefootnotemark}\,}
% \newcolumntype{Y}{>{\centering\arraybackslash}X}
% \usepackage[a4paper]{geometry}
% \usepackage{caption}
% \usepackage[bottom,flushmargin,hang,multiple]{footmisc}
% \usepackage{pdflscape}
% \usepackage[paper=portrait,pagesize]{typearea}
% \usepackage{rotating}
% \usepackage{fullpage}
% \usepackage{tabu}
\usepackage{lscape}
\newcommand{\blandscape}{\begin{landscape}}
\newcommand{\elandscape}{\end{landscape}}
\usepackage{rotating}
\newcommand{\bsideways}{\begin{sidewaystable}[htbp]}
\newcommand{\esideways}{\end{sidewaystable}}
\KOMAoption{captions}{tableheading}
\makeatletter
\@ifpackageloaded{caption}{}{\usepackage{caption}}
\AtBeginDocument{%
\ifdefined\contentsname
  \renewcommand*\contentsname{Table of contents}
\else
  \newcommand\contentsname{Table of contents}
\fi
\ifdefined\listfigurename
  \renewcommand*\listfigurename{List of Figures}
\else
  \newcommand\listfigurename{List of Figures}
\fi
\ifdefined\listtablename
  \renewcommand*\listtablename{List of Tables}
\else
  \newcommand\listtablename{List of Tables}
\fi
\ifdefined\figurename
  \renewcommand*\figurename{Figure}
\else
  \newcommand\figurename{Figure}
\fi
\ifdefined\tablename
  \renewcommand*\tablename{Table}
\else
  \newcommand\tablename{Table}
\fi
}
\@ifpackageloaded{float}{}{\usepackage{float}}
\floatstyle{ruled}
\@ifundefined{c@chapter}{\newfloat{codelisting}{h}{lop}}{\newfloat{codelisting}{h}{lop}[chapter]}
\floatname{codelisting}{Listing}
\newcommand*\listoflistings{\listof{codelisting}{List of Listings}}
\makeatother
\makeatletter
\makeatother
\makeatletter
\@ifpackageloaded{caption}{}{\usepackage{caption}}
\@ifpackageloaded{subcaption}{}{\usepackage{subcaption}}
\makeatother
\ifLuaTeX
  \usepackage{selnolig}  % disable illegal ligatures
\fi
\usepackage[round]{natbib}
\bibliographystyle{apalike}
\usepackage{bookmark}

\IfFileExists{xurl.sty}{\usepackage{xurl}}{} % add URL line breaks if available
\urlstyle{same} % disable monospaced font for URLs
\hypersetup{
  colorlinks=true,
  linkcolor={blue!000!black},
  filecolor={Maroon},
  citecolor={blue!000!black},
  urlcolor={blue!000!black},
  pdfcreator={LaTeX via pandoc}}

\author{}
\date{}

\begin{document}

\title{Regulation and bank lending in South Africa: A narrative index approach}





\author {Xolani Sibande\footnote{Economic Research Department, South African Reserve Bank; Email: xolani.sibande@resbank.co.za.} \and
Dumakude Nxumalo\footnote{Department of Economics, University of Pretoria, Pretoria, South Africa; Email: dumakude.nxumalo@up.ac.za} \and
Keaoleboga Mncube\footnote{Department of Economics, University of Pretoria, Pretoria, South Africa; Email: keaolebogamncube@gmail.com} \and
Steve Koch\footnote{Department of Economics, University of Pretoria, Pretoria, South Africa; Email: steve.koch@up.ac.za} \and
Nicola Viegi\footnote{Department of Economics, University of Pretoria, Pretoria, South Africa; Email: viegin@gmail.com}}


\date{\today}
\maketitle

\begin{abstract}
Duma.

\end{abstract}

\noindent\textbf{Keywords}: Bank lending, Narrative methods, Finance regulation\\
\textbf{JEL Codes}: G01, G18, G28, G32, G38
\newpage

\setstretch{1.5}
\section{Introduction}\label{introduction}

Inflation is known to have adverse welfare effects through its
distortion of relative prices which causes the misallocation of
resources, production inefficiencies leading to lower real economic
activity. This is well supported by both theoretical and empirical
studies. This paper contributes to the empirical literature, by
examining the relationship between inflation and relative price
variability (RPV) in South Africa given the little evidence available
and to better understand how this relationship has evolved and the
implications for monetary policy.

The South African context can provide further insights into the
relationship between inflation and RPV given the shifts in the inflation
targets from 3- 6 per cent in 2000 to 4.5 per cent since 2017 as
structural changes in inflation can affect the link between inflation
and relative price variability (Caglayan and Filiztekin (2005) and that
the functional form of relationship is U-shaped and dynamic (Fielding
and Mizen, (2000); Choi (2010)). The latter motivates the estimation of
the inflation rate that minimizes relative price distortions since if
the relationship was linear then the optimal inflation rate would be
zero. Given the inflation target reform towards a 3\% in South Africa,
this study contributes to the policy discussion on the appropriate
inflation target and understating how dynamic relations between
inflation and relative price variability has evolved.

The theoretical literature generally distinguishes between menu-cost and
signal extraction models in explaining the positive relation between
trend inflation and RVP. In these models' firms follow a one-sided
dynamic pricing rule when setting prices and face fixed non-zero costs
of adjusting nominal prices (e.g.~printing new menus) when inflation
increases as the latter reduces the optimal real price of goods and
services. Due to differences in administrative costs, these nominal
price adjustment costs vary across firms moreover these adjustments
occur at different times which implies staggard price setting behaviour
which leads to relative price variability. Menu cost models predict a
that higher expected inflation causes an increase in RVP (citations xxx)
while Lucas-type signal extraction models (Lucas(1973); Hercowitz
(1981); Cukierman (1984) xxxx more citatiosn) argue unexpected inflation
uncertainty causes higher RVP because informational frictions make it
difficult to distinguish between idiosyncratic relative shocks or an
aggregate (inflationary) shock. Menu cost and signal-extraction models
assume a positive and linear relationship making the optimal inflation
rate zero, however, recently monetary search models (Head and Kumar
(2005)) show that this relationship is nonlinear making the optimal
inflation rate somewhere between zero and positive inflation rate.

\section{Literature review}\label{literature-review}

\section{Data and methodology}\label{data-and-methodology}

\section{Results}\label{results}

\section{Conclusion}\label{conclusion}

\newpage

\section*{References}\label{references}
\addcontentsline{toc}{section}{References}

\renewcommand{\bibsection}{}
\bibliography{references.bib}

\setcounter{section}{0}
\renewcommand{\thesection}{\Alph{section}}

\setcounter{table}{0}
\renewcommand{\thetable}{A\arabic{table}}

\setcounter{figure}{0}
\renewcommand{\thefigure}{A\arabic{figure}}

\newpage

\section{Appendix}\label{appendix}

\subsection{Data}\label{data}

\begin{verbatim}
[1] FALSE  TRUE  TRUE
\end{verbatim}

\begin{verbatim}
[1] TRUE
\end{verbatim}




\end{document}
